\begin{definition}[Cumulants]
    If they exist, the ``cumulants'' \(\cumulant_n[X]\) are the coefficients of the Taylor
    series of the cumulant generating function
    \[
        \CGF_X(t)
        = \sum_{n=1}^\infty \cumulant_n[X]\frac{t^n}{n!}
        = \sum_{n=1}^\infty \CGF_X^{(n)}(0) \frac{t^n}{n!}.
    \]
\end{definition}

\begin{exercise}[subtitle={Properties of cumulants \Coffeecup}]
    Prove that cumulants have the following properties: 
    \begin{enumerate}[label=(\alph*), noitemsep]
        \item\label{it: (cumulants) accumulate} \textbf{Accumulation.}
        If \(X\) and \(Y\) are independent, then
        the cumulants ``accumulate''\footnote{
            This is where the name ``cumulant'' comes from. Since
            \(\cumulant_2[X] = \var{X}\) by Exercise \ref{ex: properties cumulant
            generating function} Bienaymé's identity (Lemma \ref{lem: Bienaymé})
            is a special case of this property for the second order cumulants.
        }
        \[
            \cumulant_n[X+Y] = \cumulant_n[X] + \cumulant_n[Y].
        \]

        \item\label{it: (cumulants) homogeneity}\textbf{Homogeneity.} For any \(a \in \real\) we have
        \[
            \cumulant_n[aX] = a^n \cumulant_n[X].
        \]

        \item\label{it: (cumulants) constant} Prove that the cumulant generating function of a constant random variable \(Y=\mu\)
        is given by
        \[
            \CGF_Y(t) = \mu t.
        \]
        For a random variable \(X\) with \(\E{X} = \mu\) deduce that
        \(\cumulant_k[X-\mu] = \cumulant_k[X]\) for all \(k \ge 2\) and
        \(\cumulant_1[X-\mu]=0\).
    \end{enumerate}
\end{exercise}

\begin{solution}
\begin{enumerate}[wide=0pt]
    \item[\ref{it: (cumulants) accumulate}]
    We have by Exercise \ref{ex: properties cumulant generating function} \ref{it: cumulant independent sums}
    \begin{align*}
        \sum_{n=1}^\infty \cumulant_n[X+Y] \frac{t^n}{n!}
        &= \CGF_{X+Y}(t)
        \\
        &= \CGF_X(t) + \CGF_Y(t)
        = \sum_{n=1}^\infty \cumulant_n[X] \frac{t^n}{n!} + \sum_{n=1}^\infty \cumulant_n[Y] \frac{t^n}{n!}.
        \\
        &= \sum_{n=1}^\infty \bigl(\cumulant_n[X] + \cumulant_n[Y]\bigr) \frac{t^n}{n!}.
    \end{align*}
    A comparison of the coefficients proves the claim.

    \item[\ref{it: (cumulants) homogeneity}]
    We have by Exercise \ref{ex: properties cumulant generating function} \ref{it: cgf linear}
    \[
        \sum_{n=1}^\infty \cumulant_n[aX] \frac{t^n}{n!}
        = \CGF_{aX}(t)
        = \CGF_X(at)
        = \sum_{n=1}^\infty \cumulant_n[X] \frac{(at)^n}{n!}
        = \sum_{n=1}^\infty a^n \cumulant_n[X] \frac{t^n}{n!}
    \]
    and the claim follows again by comparing the coefficients.

    \item[\ref{it: (cumulants) constant}] By Exercise \ref{ex: compute mgf} we have
    for a constant random variable \(Y=\mu\) that
    \[
        \MGF_Y(t) = e^{t \mu} \implies \CGF_Y(t) = t \mu.
    \]
    Comparing coefficients with the series expansion we consequently have
    \(\cumulant_1[Y] = \mu\) and \(\cumulant_k[Y] = 0\) for \(k \ge 2\).
    Since a constant random variable is independent of every other random variable,
    we have by the accumulation property that
    \[
        \cumulant_k[X-\mu]
        = \cumulant_k[X] + \cumulant_k[-\mu]
        = \cumulant_k[X]
        \quad\text{for } k \ge 2,
    \]
    and \(\cumulant_1[X-\mu] = \cumulant_1[X] + \cumulant_1[-\mu] = 0\).
    \qedhere
    \end{enumerate}
\end{solution}

\begin{remark}[Cumulants of normal distribution]
    Recall that the cumulant generating function of the 
    normal distribution is a quadratic polynomial
    \[
        \CGF_X(t) = \mu t + \frac12 \sigma^2 t^2.
    \]
    Thus all cumulants of order three and higher are zero! It is the only
    distribution with this property. One can show that the cumulants of order
    three and higher vanish for \(\sqrt{n}\, \mean{X}_n\) as \(n\to\infty\)
    (Exercise \ref{ex: cumulant CLT}). This can be turned into a weak
    form of the central limit theorem, which requires that the moment generating
    function (and therefore all moments) are finite.
\end{remark}

\begin{exercise}[subtitle={Cumulant CLT}]
    \label{ex: cumulant CLT}
    Let \(X_1, X_2, \ldots\) be independent copies of a random variable
    \(X\) with mean \(\mu\), variance \(\sigma^2\) and finite moments of all orders.
    Let \(\mean{X}_n = \frac1n \sum_{i=1}^n X_i\) be the empirical mean.
    Show that the cumulants \(\cumulant_k[\sqrt{n}(\mean{X}_n - \mu)]\)
    of the random variable \(\sqrt{n}(\mean{X}_n - \mu)\) satisfy
    \begin{align*}
        \cumulant_1[\sqrt{n}(\mean{X}_n - \mu)] &= 0,
        \\
        \cumulant_2[\sqrt{n}(\mean{X}_n - \mu)] &= \sigma^2,
        \\
        \cumulant_k[\sqrt{n}(\mean{X}_n - \mu)] &\to 0
        \quad\text{for } k \ge 3.
    \end{align*}
    Conclude that the cumulant generating function of
    \(\sqrt{n}(\mean{X}_n - \mu)\) converges to the cumulant generating function
    of the normal distribution \(\normal{0, \sigma^2}\).
\end{exercise}

\begin{solution}
    By Exercise \ref{ex: properties cumulant generating function}
    the first and second cumulant are the expectation and variance
    respectively and therefore
    \begin{align*}
        \cumulant_1[\sqrt{n}(\mean{X}_n - \mu)]
        &= \E{\sqrt{n}(\mean{X}_n - \mu)} = 0,
        \\
        \cumulant_2[\sqrt{n}(\mean{X}_n - \mu)]
        &= \var{\sqrt{n}(\mean{X}_n - \mu)}
        = n \var{\mean{X}_n}
        = n \tfrac{\sigma^2}{n} = \sigma^2.
    \end{align*}
    For the higher order cumulants we have we use the accumulation property
    and homogeneity of Exercise \ref{ex: properties cumulant generating function} to get
    \begin{align*}
        \cumulant_k[\sqrt{n}(\mean{X}_n - \mu)]
        &= \cumulant_k\Bigl[\frac1{\sqrt{n}} \sum_{i=1}^n (X_i - \mu)\Bigr]
        \\
        \overset{\text{accum.}}&= n\, \cumulant_k[\tfrac1{\sqrt{n}} (X_1 - \mu)]
        \\
        \overset{\text{hom.}}&= n^{1-\frac{k}{2}} \cumulant_k[X_1 - \mu].
        \\
        &= n^{1-\frac{k}{2}} \cumulant_k[X_1] \qquad \forall k \ge 2.
    \end{align*}
    Clearly for \(k\ge 3\) the cumulants decay to zero. Convergence of the
    cumulant generating function then follows by dominated convergence since the
    series of analytic functions converges absolutely in a neighborhood of zero.
\end{solution}
